% -------------------------
%
% -NOTACIÓN: Lista de símbolos, acrónimos y glosario de términos.
%
% -------------------------
\cleardoublepage
\phantomsection % Necesario con hyperref

\chapter*{Notación y acrónimos} % Opción con * para que no aparezca en TOC ni numerada
\addcontentsline{toc}{chapter}{Notación y acrónimos} % Añade al TOC.

\section*{Notacion}

\begin{tabular}{r r p{0.78\linewidth}}
$z_c$        & : & \emph{Logit} del código $c$: salida del clasificador antes de la sigmoide. Puede tomar cualquier valor real \\
$\sigma(z)$  & : & Función sigmoide, que convierte un \emph{logit} en una probabilidad entre 0 y 1 \\
$\tau$       & : & Umbral de decisión: un código se propone si $\sigma(z_c) \geq \tau$ \\
$\beta$      & : & Peso del diccionario en el modo fusionado: $s_c = z_c + \beta \cdot \mathrm{conf}(\mathrm{bloque}(c))$ \\
$w_c$        & : & \texttt{pos\_weight} de la clase $c$: peso del término positivo en la pérdida, que compensa el desequilibrio \\
$|W|$        & : & Número de palabras candidatas de un informe en el cálculo de la atribución \\
\end{tabular}

\section*{Lista de acrónimos}
% OJO: Esta lista debería estar ordenada alfabeticamente (hacer de modo manual).

\begin{tabular}{r r p{0.78\linewidth}}
API   & : & \emph{Application Programming Interface}, interfaz de programación \\
BCE   & : & \emph{Binary Cross-Entropy}, entropía cruzada binaria; función de pérdida del clasificador \\
BERT  & : & \emph{Bidirectional Encoder Representations from Transformers} \\
BSC   & : & Barcelona Supercomputing Center \\
CIE-10 & : & Clasificación Internacional de Enfermedades, décima revisión \\
CIE-10-ES & : & Edición española de la CIE-10, la que usa el Sistema Nacional de Salud \\
CORS  & : & \emph{Cross-Origin Resource Sharing}, control de peticiones entre dominios \\
CPU   & : & \emph{Central Processing Unit}, procesador de propósito general \\
CQRS  & : & \emph{Command Query Responsibility Segregation}, separación entre el modelo que escribe y el que lee \\
CSV   & : & \emph{Comma-Separated Values}, formato tabular en texto plano \\
DSN   & : & \emph{Data Source Name}, cadena de conexión; aquí identifica el proyecto de Sentry \\
F1    & : & Media armónica de precisión y \emph{recall} \\
GGUF  & : & Formato de pesos cuantizados que emplea \texttt{llama.cpp} \\
GPU   & : & \emph{Graphics Processing Unit}, acelerador empleado en entrenamiento e inferencia \\
HSTS  & : & \emph{HTTP Strict Transport Security}, cabecera que fuerza el uso de HTTPS \\
HTTP  & : & \emph{HyperText Transfer Protocol} \\
LLM   & : & \emph{Large Language Model}, modelo de lenguaje generativo \\
MAP   & : & \emph{Mean Average Precision}, métrica oficial del \emph{benchmark} CodiEsp (véase el glosario) \\
NDJSON & : & \emph{Newline-Delimited JSON}, un objeto JSON por línea; formato del resumen en \emph{streaming} \\
OTP   & : & \emph{Open Telecom Platform}, biblioteca de concurrencia y supervisión de Erlang \\
PLN   & : & Procesamiento del Lenguaje Natural (en inglés, NLP) \\
REST  & : & \emph{Representational State Transfer}, estilo arquitectónico de la API \\
RGPD  & : & Reglamento General de Protección de Datos, Reglamento (UE) 2016/679 \\
SDK   & : & \emph{Software Development Kit}, biblioteca de integración \\
SMTP  & : & \emph{Simple Mail Transfer Protocol}, protocolo de envío de correo \\
TLS   & : & \emph{Transport Layer Security}, cifrado del tráfico en tránsito \\
UUID  & : & \emph{Universally Unique Identifier} \\
ZLPR  & : & \emph{Zero-bounded Log-sum-exp \& Pairwise Rank-based loss}, pérdida de ordenación multietiqueta \\
\end{tabular}

\section*{Glosario de términos}

Términos que el texto emplea con un significado preciso y que conviene fijar antes de leerlo.

\begin{tabular}{r r p{0.78\linewidth}}
Atribución & : & Palabras del informe que sostienen un código predicho. Se calcula enmascarándolas y midiendo cuánto cae el \emph{logit} \\
\emph{Benchmark} & : & Tarea compartida con datos y métrica fijos que permite comparar sistemas. Aquí, CodiEsp-D 2020 \\
Bloque & : & Los tres primeros caracteres de un código CIE-10. Agrupa códigos de la misma familia (por ejemplo, \texttt{E11} para la diabetes tipo 2) \\
Codificador clínico & : & La persona que asigna los códigos a un informe. No confundir con el \emph{encoder} del modelo \\
Cola larga & : & Distribución en la que unos pocos códigos concentran casi todos los ejemplos y la mayoría aparece muy pocas veces \\
\emph{Encoder} & : & La parte del modelo que convierte el texto en vectores. Aquí, RigoBERTa-Clinical \\
\emph{Event Sourcing} & : & Guardar la secuencia inmutable de hechos ocurridos en lugar del estado final, que se deriva de ellos \\
\emph{Fine-tuning} & : & Ajustar un modelo ya preentrenado sobre la tarea concreta, en vez de entrenarlo desde cero \\
\emph{Logit} & : & Salida del clasificador antes de aplicar la sigmoide. Ordena los códigos aunque no sea una probabilidad \\
MAP & : & Para cada informe se mide la precisión media de su lista ordenada de códigos, y se promedian esos valores entre informes. No depende del umbral: mide el orden, no la decisión \\
MAP alcanzable & : & MAP calculado solo sobre los códigos que el modelo puede predecir. No es comparable con el de otros sistemas \\
MAP estricto & : & MAP sobre el \emph{gold} completo, penalizando los códigos fuera del vocabulario. Es el comparable con el \emph{benchmark} \\
Modo fusionado & : & Motor que suma la confianza del diccionario al \emph{logit} del clasificador y ordena con la puntuación resultante \\
Multietiqueta & : & Cada informe puede llevar varios códigos a la vez, de modo que las clases no se excluyen entre sí \\
\emph{Pipeline} & : & Cadena de etapas de procesado encadenadas \\
\texttt{pos\_weight} & : & Peso que la pérdida da a los ejemplos positivos de cada clase, para que los códigos raros no se ignoren \\
Proyección & : & Tabla de solo lectura derivada de los eventos, que es la que consulta la aplicación \\
\emph{Ranking} & : & Lista de códigos ordenada por confianza, sin decidir cuáles se aceptan \\
\emph{Recall} & : & Proporción de los códigos correctos que el sistema llega a proponer \\
Retrotraducción & : & Generar variantes de un texto traduciéndolo a otra lengua y de vuelta, para aumentar el corpus \\
Umbral & : & Valor a partir del cual una probabilidad se convierte en un código propuesto. Afecta al F1, no al MAP \\
Vocabulario del modelo & : & Los 1.767 códigos presentes en el entrenamiento, que son los únicos que el clasificador puede proponer \\
\end{tabular}
